\documentclass{article}

\usepackage[utf8]{inputenc}
\usepackage[T1]{fontenc}
\usepackage{helvet}
\usepackage{geometry}
\usepackage{xcolor}
\usepackage{eso-pic}
\usepackage{fancyhdr}
\usepackage{parskip}
\usepackage{microtype}
\usepackage[english, ukrainian]{babel}
\usepackage{url}
\usepackage{amsmath}
\usepackage{amssymb}
\usepackage{amsthm}
\usepackage{longtable}
\usepackage{booktabs}
\usepackage{hyperref}
\usepackage{titlesec}

\definecolor{nato}{HTML}{293757}
\definecolor{footergray}{gray}{0.65}

\geometry{a4paper,left=3cm,right=3cm,top=3cm,bottom=3cm}
\renewcommand{\familydefault}{\sfdefault}
\setcounter{secnumdepth}{3}

\titleformat{\section}
  {\color{nato}\Huge\bfseries}{\thesection.}{1em}{}
\titlespacing*{\section}{0pt}{30pt}{12pt}

\titleformat{\subsection}
  {\color{nato}\LARGE\bfseries}{\thesubsection.}{1em}{}
\titlespacing*{\subsection}{0pt}{20pt}{8pt}

\titleformat{\subsubsection}
  {\color{nato}\Large\bfseries}{\thesubsubsection.}{1em}{}
\titlespacing*{\subsubsection}{0pt}{10pt}{6pt}

\fancyhf{}
\fancyhead[R]{\small\color{footergray} 27 червня 2026}
\fancyfoot[L]{\small\color{footergray} Технічні вимоги ERP/1. Модуль Документи}
\fancyfoot[R]{\small\color{footergray}\thepage}

\newcommand\HUGE[1]{\fontsize{40}{40}\selectfont #1}

\renewcommand{\headrulewidth}{0pt}
\renewcommand{\footrulewidth}{0.4pt}
\renewcommand{\footrule}{\color{footergray}\hrule width \textwidth height 0.4pt}

\begin{document}

\pagestyle{fancy}

\begin{titlepage}
  \AddToShipoutPictureBG*{\AtPageLowerLeft{\color{nato}\rule{\paperwidth}{\paperheight}}}
  \color{white}\thispagestyle{empty}\vspace*{3cm}\centering
  {\Huge  \textbf{ERP/1: Документи} \par} \vspace{1.5cm}
  {\HUGE  \textbf{Технічні вимоги} \par} \vspace{1.5cm}
  {\LARGE \textbf{до системи електронного документообігу \\ та автоматизації діловодства \\ в установах і міжвідомчого обміну} \par} \vspace{1.5cm}
  {\large Адаптовано на основі N2O.DEV \\ та нормативно-технічних вимог МВС України \par}
  \vfill
  {\large 2026 \copyright\ Системи Електронного Урядування \par}
\end{titlepage}

\newpage

\begin{center}
  {\Huge\color{nato}\bfseries Зміст\par}
\end{center}
\vspace{40pt}
\tableofcontents

\newpage

\section{Умовні скорочення та визначення}

\begin{longtable}{p{3.5cm}p{10.5cm}}
\toprule
\textbf{Термін / Абревіатура} & \textbf{Визначення} \\
\midrule
СЕД & Система електронного документообігу. \\
СЕВ ОВВ & Система електронної взаємодії органів виконавчої влади. \\
КЕП & Кваліфікований електронний підпис (X.509 CAdES-X Long). \\
КЕПч & Кваліфікована електронна печатка установи. \\
РМК & Реєстраційно-моніторингова картка документа. \\
ОШС & Організаційно-штатна структура установи. \\
НДІ & Нормативно-довідкова інформація (довідники, класифікатори). \\
BPE & Business Process Engine (асинхронне ядро бізнес-процесів). \\
N2O & Протокол та веб-фреймворк для реактивної взаємодії через WebSockets. \\
ABAC & Attribute-Based Access Control (контроль доступу на основі атрибутів). \\
MHS X.420 & Message Handling System (протокол поштових конвертів ISO 10021-6). \\
X.422.3 & Протокол телекомунікаційного обміну службовими повідомленнями СЕВ. \\
\bottomrule
\end{longtable}

\section{Загальні відомості}

\subsection{Передумови розробки}
Створення підсистеми «ERP/1: Документи» обумовлене необхідністю переходу державних органів та установ України на високопродуктивну, безпечну та відмовостійку хмарну веб-платформу. Система повністю покриває вимоги Постанови КМУ № 55 щодо документування управлінської діяльності та організації міжвідомчого документообігу на Erlang/OTP технологіях.

\subsection{Законодавча та нормативна основа}
Підсистема розробляється та функціонує відповідно до таких законодавчих та підзаконних нормативно-правових актів України:

\subsubsection{Загальне законодавство України}
\begin{itemize}
    \item Закон України «Про інформацію» (№ 2657-XII);
    \item Закон України «Про Національну програму інформатизації» (№ 74/98-ВР);
    \item Закон України «Про електронні документи та електронний документообіг» (№ 851-IV);
    \item Закон України «Про електронні довірчі послуги» (№ 2155-VIII);
    \item Закон України «Про звернення громадян» (№ 393/96-ВР);
    \item Закон України «Про доступ до публічної інформації» (№ 2939-VI);
    \item Розпорядження КМУ від 15 травня 2013 р. № 386-р «Про схвалення Стратегії розвитку інформаційного суспільства в Україні»;
    \item Постанова КМУ від 26 травня 2006 р. № 373 «Про затвердження Правил забезпечення захисту інформації в інформаційних, телекомунікаційних та інформаційно-телекомунікаційних системах».
\end{itemize}

\subsubsection{Законодавство у сфері електронного документообігу та КЕП}
\begin{itemize}
    \item Постанова Кабінету Міністрів України від 17 січня 2018 р. № 55 «Деякі питання документування управлінської діяльності»;
    \item Наказ Міністерства юстиції України від 1600/5 «Про затвердження Порядку роботи з електронними документами через систему електронної взаємодії органів виконавчої влади з використанням електронного цифрового підпису»;
    \item Постанова Кабінету Міністрів України від 7 листопада 2018 р. № 749 «Про затвердження Порядку використання електронних довірчих послуг в органах державної влади, органах місцевого самоврядування, підприємствах, установах та організаціях державної форми власності»;
    \item Наказ Міністерства цифрової трансформації України та Адміністрації ДССЗЗІ № 140/614 щодо експертизи засобів захисту;
    \item Наказ Міністерства освіти і науки, молоді та спорту України від 1207 «Про вимоги до форматів даних електронного документообігу в органах державної влади. Формат електронного повідомлення»;
    \item Наказ Міністерства юстиції України від 1886/5 «Про затвердження Порядку роботи з електронними документами у діловодстві та їх підготовки до передавання на архівне зберігання»;
    \item Наказ ДП «УкрНДНЦ» від № 144 про прийняття та скасування національних стандартів ДСТУ 4163:2020 та ДСТУ 9031:2020.
\end{itemize}

\subsubsection{Законодавство для спеціалізованих додаткових модулів}
\begin{itemize}
    \item \textbf{Модуль Провадження (Прокуратура/Суд):} Кримінальний процесуальний кодекс України (№ 4651-VI); Наказ Офісу Генерального прокурора № 298 «Про затвердження Положення про Єдиний реєстр досудових розслідувань, порядок його формування та ведення».
    \item \textbf{Модуль Закупівлі:} Закон України «Про публічні закупівлі» (№ 922-19); Постанова КМУ № 169 та Постанова КМУ № 1178; Закон України «Про оборонні закупівлі» (№ 808-20); Постанови КМУ № 1275, № 1070, № 224, № 710, № 544, № 1495; Накази Мінекономіки № 708 та № 1749.
    \item \textbf{Модуль Зброя:} Проекти Законів про право на цивільну вогнепальну зброю (№ 5708, № 5709) та інструкції МВС щодо Єдиного реєстру зброї.
\end{itemize}

\section{Призначення та цілі впровадження}
Модуль «Документи» призначений для організації проходження та автоматизації життєвого циклу електронних документів, візування, підписання КЕП, контролю за виконанням завдань, а також міжвідомчого поштового обміну за стандартами X.420 / X.422.3.
\\
\textbf{Цілі створення засобу:}
\begin{itemize}
    \item Скорочення часу обробки та реєстрації вхідних, вихідних та внутрішніх документів;
    \item Забезпечення прозорості та повного аудиту дій над документами (Audit Log);
    \item Економія матеріальних ресурсів установи на паперовий обіг;
    \item Забезпечення юридичної сили за допомогою накладання КЕП та КЕПч;
    \item Оперативний пошук та побудова аналітичної звітності виконавської дисципліни.
\end{itemize}

\section{Класифікація вимог}

\subsection{Функціональні вимоги}

\subsubsection{Словникова та довідникова система}
Платформа повинна підтримувати та використовувати ієрархічні класифікатори (Дерева) та словники:
\begin{itemize}
    \item \textbf{ІНСТРУКЦІЯ:} Довідник нормативно-методологічних правил діловодства.
    \item \textbf{КОАТУУ:} Класифікатор об'єктів адміністративно-територіального устрою України.
    \item \textbf{ПРОЦЕСИ:} Довідник бізнес-процесів та маршрутів документів.
    \item \textbf{СЛОВНИКИ:} Загальні класифікатори (наприклад, типи звернень громадян, види контролю).
\end{itemize}

\subsubsection{Опис реєстрових сутностей (ERP-Схема даних)}
Доменна модель системи повинна включати такі сутності:
\begin{itemize}
    \item \textbf{ORGANIZATION:} Установа, контрагент (ЄДРПОУ, юридична адреса).
    \item \textbf{BRANCH:} Філія, територіальний орган чи структурний підрозділ.
    \item \textbf{EMPLOYEE:} Співробітник (посада, підрозділ, контактні дані).
    \item \textbf{PERSON:} Фізична особа (ПІБ, ІПН, сертифікати КЕП).
    \item \textbf{SIGNATURE:} Реквізити електронного підпису (підписувач, сертифікат, TSP мітка часу).
    \item \textbf{ABAC:} Об'єкт прав доступу (суб'єкт, дія, контекстні обмеження).
    \item \textbf{FILEDESC:} Специфікація медіа-файлу (хеш SHA-256, назва, розмір, RocksDB посилання).
    \item \textbf{PROJECT:} Картка проекту документа на етапі погодження.
    \item \textbf{POSITION:} Штатна посада користувача.
    \item \textbf{ADDRESS:} Поштові та фізичні адреси установ та контрагентів.
    \item \textbf{LOCATION:} Географічні координати.
    \item \textbf{DEED:} Справа архівного зберігання (номенклатурна папка).
    \item \textbf{CATDOC:} Категорії та види документів.
    \item \textbf{CATDEED:} Класифікатор справ.
    \item \textbf{NOMDEED:} Номенклатура справ на рік.
\end{itemize}

\subsubsection{Опис бізнес-процесів (BPE Flows)}
Усі процеси моделюються як FSM автомати в ядрі BPE:
\begin{itemize}
    \item \textbf{INPUT:} Обробка вхідних документів від реєстрації до виконання.
    \item \textbf{OUTPUT:} Створення, візування, підписання КЕП та відправка вихідного листа.
    \item \textbf{INTERNAL:} Обробка внутрішніх документів (службові записки, накази).
    \item \textbf{ORG:} Підготовка та доведення організаційно-розпорядчих документів.
    \item \textbf{CONTROLTASK:} Контроль виконання доручень.
    \item \textbf{EXECTASK:} Виконання завдань співвиконавцями.
    \item \textbf{ADMISSION:} Реєстрація та контроль розгляду звернень громадян.
    \item \textbf{ADDITIONAL:} Додаткові асинхронні бізнес-процеси.
\end{itemize}

\subsubsection{Опис типів документів}
Система повинна опрацьовувати такі типи документів (категорії РМК):
\begin{itemize}
    \item \textbf{INPUT:} Вхідний лист від установи чи фізичної особи.
    \item \textbf{OUTPUT:} Вихідний лист або відповідь на звернення.
    \item \textbf{INTERNAL:} Внутрішній документ (заява, доповідна записка).
    \item \textbf{ORG:} Наказ керівника установи з основної діяльності.
    \item \textbf{CITIZEN:} Звернення громадянина відповідно до ЗУ «Про звернення громадян».
    \item \textbf{ESQUIRE:} Адвокатський запит (скорочений термін розгляду до 5 робочих днів).
    \item \textbf{PUBLIC:} Запит на публічну інформацію (термін розгляду до 5 робочих днів).
    \item \textbf{TASK \& SUBTASK:} Доручення (резолюція) та підлеглі підзадачі виконавцям.
    \item \textbf{ADDITIONAL:} Додаткові супровідні документи.
\end{itemize}

\subsubsection{Опис сторінок інтерфейсу клієнта (Controllers)}
Клієнтський SPA-додаток на базі N2O/Nitro повинен містити такі сторінки:
\begin{itemize}
    \item \textbf{LDAP:} Дерево підрозділів, пошук та перегляд профілів співробітників.
    \item \textbf{HELP:} Контекстна довідкова система для користувачів.
    \item \textbf{KVS:} Інспектор сховища RocksDB/Mnesia для адміністраторів системи.
    \item \textbf{BPE:} Панель керування активними FSM процесами.
    \item \textbf{CRM:} Головна робоча область із реєстрами документів.
    \item \textbf{RMK:} Екран перегляду та редагування РМК конкретного документа.
    \item \textbf{ROLES:} Управління правами доступу користувачів.
    \item \textbf{ROUTES:} Маршрутна карта переходів інтерфейсу.
\end{itemize}

\subsubsection{Елементи інтерфейсу та Комболукапи}
Для швидкого автозаповнення та рендерингу складних форм використовуються Nitro-компоненти:
\begin{itemize}
    \item \textbf{Комболукапи:} ADDRESS, DEPUTY, DICT, EMPLOYEE, BRANCH, SEARCH, SEV, ROLE, ORG, LOCATION, PLACE, POSITION.
    \item \textbf{Елементи інтерфейсу:} ABAC (права), BPE (процеси), CRM (реєстр), SEARCH (динамічний пошук), UPLOAD (завантаження файлів), ATTACH (асоціація додатків), RMK (картка), DOCS (документи), FILES (файли), MONITOR (моніторинг), ROLE (ролі), TABLE (таблиця), TREE (дерева).
    \item \textbf{Міні-редактори:} DICT, EMPLOYEE, PLACE, BRANCH, ORG, FILE, EDIT, MONITOR, PROCESS, TRACE.
    \item \textbf{Автентифікаційні форми:} ACTIVATION (активація КЕП), KEY (вибір носія), FILE (завантаження ключа), ANY (інші методи), PROFILE (профіль), ERROR (помилки).
\end{itemize}

\subsubsection{Шаблони документів}
Автоматична генерація файлів за затвердженими бланками установи:
\begin{itemize}
    \item \textbf{ЛИСТ:} Шаблон офіційного вихідного листа.
    \item \textbf{БЛАНК:} Загальний бланк установи для наказів та рішень колегії.
    \item \textbf{ЗАПИСКА:} Шаблон внутрішньої службової записки.
    \item \textbf{ВІДПОВІДЬ:} Шаблон відповіді на звернення чи запит.
\end{itemize}

\subsubsection{Галузеві додаткові модулі}
\begin{itemize}
    \item \textbf{МІА: Провадження (Судочинство):} облік судових та кримінальних проваджень, клопотань та ухвал суду.
    \item \textbf{МІА: Закупівлі:} облік публічних закупівель (договори, процедури за Постановами КМУ № 169/1178).
    \item \textbf{МІА: Зброя:} інтеграція з Єдиним реєстром зброї, облік дозволів та власників.
\end{itemize}

\subsubsection{СЕВ ОВВ Інтеграція}
Обмін документами між державними органами через СЕВ ОВВ здійснюється за допомогою:
\begin{itemize}
    \item \textbf{ДОКУМЕНТ:} Реквізити обміну електронного документа.
    \item \textbf{НПА:} Реквізити нормативно-правового акту.
    \item \textbf{ЗАДАЧА:} Синхронізація завдань та звітів виконання між відомствами.
    \item Передача через поштові конверти \textbf{MHS X.420} та транспортний протокол \textbf{X.422.3}.
\end{itemize}

\subsubsection{Автентифікація та КЕП}
Вхід здійснюється за допомогою X.509 сертифікатів КЕП. Підписання документів виконується у форматі CAdES-X Long (із вбудованою позначкою часу TSP).

\subsection{Нефункціональні вимоги}

\subsubsection{Вимоги до архітектури}
\begin{itemize}
    \item \textbf{Модульність:} Система повинна містити чітко розмежовані сервіси та модулі (близько 100 модулів).
    \item \textbf{База даних:} Native Erlang Mnesia + RocksDB (без реляційного SQL) для максимальної швидкості транзакцій.
    \item \textbf{Клієнт-Сервер:} WebSocket транспорт N2O з Data-on-Wire JSON передачею.
\end{itemize}

\subsubsection{Вимоги до безпеки}
\begin{itemize}
    \item Накладання КЕПч установи на всі файли, що надсилаються через СЕВ ОВВ.
    \item Перевірка сертифікатів через OCSP/CRL.
    \item Захист сесій через Secure Cookie та перевірка ABAC правил доступу.
\end{itemize}

\subsubsection{Вимоги до продуктивності}
\begin{itemize}
    \item Рендеринг РМК --- до 500 мс.
    \item Підтримка до 10,000 користувачів одночасно.
\end{itemize}

\subsubsection{Вимоги до логування}
\begin{itemize}
    \item Кожна дія над документом (перегляд, зміна статусу, КЕП підпис) повинна фіксуватися в Audit Log у форматі JSON.
    \item Передача логів до стеку ELK в реальному часі.
\end{itemize}

\newpage
\appendix
\section{Специфікація реквізитів сутностей}

\begin{longtable}{p{3cm}p{3cm}p{2cm}p{6cm}}
\caption{Реквізити картки вхідного документа (INPUT)} \\
\toprule
\textbf{Поле} & \textbf{Тип даних} & \textbf{Обов'язкове} & \textbf{Опис призначення} \\
\midrule
\texttt{id} & UUID & Так & Унікальний ідентифікатор РМК. \\
\texttt{doc\_number} & String & Так & Реєстраційний номер. \\
\texttt{doc\_date} & Date & Так & Дата реєстрації документа. \\
\texttt{correspondent} & UUID & Так & Посилання на організацію-відправника (ЄДРПОУ). \\
\texttt{subject} & Text & Так & Тема (короткий зміст) листа. \\
\texttt{urgent} & Boolean & Так & Ознака терміновості розгляду. \\
\texttt{due\_date} & Date & Ні & Контрольний строк виконання. \\
\texttt{files} & List(UUID) & Так & Список прикріплених файлів (FileDesc). \\
\texttt{status} & Enum & Так & Поточний стан (draft, registered, archived). \\
\bottomrule
\end{longtable}

\begin{longtable}{p{3cm}p{3cm}p{2cm}p{6cm}}
\caption{Реквізити Задачі / Резолюції (Task)} \\
\toprule
\textbf{Поле} & \textbf{Тип даних} & \textbf{Обов'язкове} & \textbf{Опис призначення} \\
\midrule
\texttt{id} & UUID & Так & Унікальний ідентифікатор задачі. \\
\texttt{rmk\_id} & UUID & Так & Посилання на РМК документа. \\
\texttt{author\_id} & UUID & Так & Керівник, який наклав резолюцію. \\
\texttt{chief\_executor} & UUID & Так & Головний відповідальний виконавець. \\
\texttt{executors} & List(UUID) & Так & Список співвиконавців. \\
\texttt{description} & Text & Так & Текст доручення. \\
\texttt{due\_date} & Date & Так & Термін виконання. \\
\texttt{status} & Enum & Так & Стан (waiting, active, completed, cancelled). \\
\bottomrule
\end{longtable}

\section{Технічні деталі реалізації процесів та валідацій}

\subsection{Erlang Records для СЕД}
Базові записи збереження даних документообігу:
\begin{verbatim}
-record(document_rmk, {
    id,
    doc_number,
    doc_date,
    doc_type,
    correspondent_id,
    subject,
    author_id,
    files = [],
    status = draft
}).

-record(document_task, {
    id,
    parent_id,
    rmk_id,
    author_id,
    executors = [],
    chief_executor,
    description,
    due_date,
    status = waiting
}).
\end{verbatim}

\subsection{Модуль валідації електронних підписів kep\_validator.erl}
Код Erlang-модуля для перевірки довіри до КЕП:
\begin{verbatim}
-module(kep_validator).
-export([validate_signature/2, verify_time_stamp/1]).

validate_signature(FileHash, SignatureBytes) ->
    case ca_client:verify_signature(FileHash, SignatureBytes) of
        {ok, SignerCert} ->
            case ca_client:check_cert_status(SignerCert) of
                valid -> ok;
                expired -> {error, cert_expired};
                revoked -> {error, cert_revoked}
            end;
        Error -> Error
    end.

verify_time_stamp(TimeStampToken) ->
    case ca_client:verify_tsp(TimeStampToken) of
        true -> ok;
        false -> {error, invalid_tsp}
    end.
\end{verbatim}

\end{document}
